\chapter{Introduction}\label{ch:introduction}

\section{Motivation}\label{sec:motivation}

The same Wikipedia article on a contested historical event reads
differently depending on which language edition one opens.
\citet{oeberst2020wikipedia} make this observation precise. Across
thirty-five intergroup conflicts and three studies, they show that
the language edition of each conflicting party portrays the
in-group as preferred, more powerful, less immoral, and less
responsible for the conflict, with the strength of bias predicted
by conflict recency and the share of in-group editors among the
article's top contributors. The unit of analysis is the
human-authored article lead --- exactly the corpus that
contemporary multilingual large language models are trained on.
The natural follow-on, taken up by this thesis, is whether the
same divergence appears when the curator is no longer a human
editor but a language-conditional LLM. Throughout, \emph{ingroup
bias} is used in this Oeberst-inherited sense --- systematic
resemblance to the in-group's own account of an event, a property
of content measurable without imputing allegiance; whether
resemblance extends to taking the in-group's side is exactly what
the stance instruments of Chapter~\ref{ch:experiments} test. One
scoping note applies from the start: where the prompt language is
not the language of a conflict party --- English on the
Russo-Ukrainian pair, for instance --- the same quantity is better
read as own-language \emph{corpus} pull than as ingroup bias in
Oeberst's party sense; \S\ref{sec:four_instruments} develops this
distinction.

The shift in actor matters. Wikipedia's edit history keeps
disagreement visible and contestable; an LLM response is a single
answer whose provenance is opaque by construction, so an inherited
bias that was open on Wikipedia becomes invisible in the model. The
scale shifts too: language editions are read by communities who
already share the language, while LLMs are queried by anyone in
their preferred language --- including readers consulting them
precisely because they cannot verify the answer in their own
language, which is exactly the contested-event case.

A precedent literature has begun to map this terrain.
\citet{durmus2023globalopinionqa} prompt frontier LLMs with
opinion questions in a respondent's native language and report
that the responses do not move toward that respondent's
national opinion distribution --- a null result on the simplest
version of linguistic prompting. \citet{buyl2024llm} take the
opposite side: when the prompt asks about a political person,
the model's implied evaluation correlates with the geopolitical
alignment of the model's creator country, particularly under
post-training alignment. \citet{urman2025silence} audit three Western
chatbots on politically sensitive queries and find that Google
Bard refuses 90\% of Russian-language Putin queries against
only 19\% of the English equivalents, with the Russian
non-response pattern matching the leaked Roskomnadzor block
list more closely than Bard's own published content policy. \citet{smirnov2026language} provides the
qualitatively closest parallel: prompting the same model with
semantically equivalent Russian and Ukrainian renderings of a
Ukrainian civil-society document produces Russian state-
discourse vocabulary in the Russian condition and Western
liberal-democratic vocabulary in the Ukrainian condition.
The model and the source content are held fixed; the prompt
language alone shifts the framing.

The contribution of this thesis is not to claim that
language-conditional bias in multilingual LLMs is a novel
discovery. It is to provide the multi-event, multi-provider,
embedding-quantified, English-controlled, and
imaginary-falsified extension that the precedents themselves
flag as missing. Where \citet{durmus2023globalopinionqa} test
one methodology on opinion questions and report a null, this
thesis tests the same hypothesis on contested-event
\emph{factual} questions across nine prompts in three languages
on the 2014 Russo-Ukrainian crisis, then scales to five
conflict families and nine languages. Where
\citet{buyl2024llm} evaluate nineteen LLMs across six UN
languages but with the four major Western and Chinese provider
ecosystems represented as the comparison axis, this thesis
sweeps fifteen providers across seven regulatory regimes --- the
United States, China (CAC-regulated), Canada (Cohere multilingual),
Israel, Saudi Arabia, Taiwan, and Russia (state-aligned) --- on the
same prompt bank. Fourteen of the fifteen return analysable prose;
the Russian provider refuses uniformly and is analysed separately as
evidence in its own right. Where \citet{urman2025silence}
document \emph{that} refusal patterns inherit from Russian
state templates, this thesis recovers the embedding-space
signature of those refusals through unsupervised clustering,
closing the explicit creator-provider gap left open by that
paper. Where \citet{smirnov2026language} establishes the
pattern qualitatively on one document, this thesis quantifies
the same shift geometrically across thousands of responses and
shows when the literal pattern survives, when it dissolves,
and when an alternative metric --- a paired-edit stance-axis
projection --- recovers a stronger signal than anchor cosine.

Three methodological commitments shape what follows. First,
English is treated as a control language rather than as one of
$N$ symmetric options. Frontier LLM training corpora are
English-dominant; the asymmetry between an English baseline
and non-English target languages is itself a measurement, and
collapsing it loses information. Second, the analysis is built
around \emph{anchors} --- the multilingual Wikipedia article
on each event in each language --- so that cross-lingual
divergence is measured in a frame that does not require
native-speaker raters. The anchor frame inherits its
central assumption from \citet{oeberst2020wikipedia}:
Wikipedia language editions encode each side's framing of a
contested event. When the anchor pre-exists in three
languages the response embedding can be projected against all
three; when the event is fictional the anchor frame collapses
and a complementary metric is required. Both regimes are
treated. Third, the design is depth-first. Every instrument in
the thesis --- anchor cosine, language-axis debiasing,
unsupervised structure recovery, the stance axis, the judge
protocols --- is built and validated on a single case, the 2014
Russo-Ukrainian crisis, where the anchor corpus is richest, the
precedent literature is concentrated, and the prompt
translations receive the closest scrutiny. The other conflict
families and the sixteen-language event-ladder corpus then
serve a different role: they test whether the validated
instruments generalise, at lower sampling depth and with
machine-translated prompts. The thesis is about multilingual
LLMs in general; the Russo-Ukrainian case is its evidentiary
spine, and the asymmetry between the two is a design choice
rather than an accident of emphasis. The choice is vindicated
in \S\ref{sec:exp_021}: the home case turns out to be atypical
on the stance instrument --- a fact only the breadth arm could
have revealed.

\section{Research questions}\label{sec:rqs}

This thesis investigates four coupled questions.

\begin{description}
  \item[\textbf{RQ1 --- Ingroup pull.}] When prompted in
    language $\ell$ about a contested event $q$, does an LLM's
    response embed closer to the $\ell$-language Wikipedia
    article on $q$ than to the other-language articles, after
    controlling for the baseline same-language similarity
    inherent in shared vocabulary?
  \item[\textbf{RQ2 --- Generality across providers and
    conflicts.}] Does any such ingroup pull survive across
    providers from distinct training-data ecosystems
    (US-trained, China CAC-regulated, Russian state-aligned,
    Cohere multilingual, Saudi state-research, Taiwan
    state-counter, Israeli) and across contested-event
    families beyond the Russo-Ukrainian case (Israel--Palestine,
    India--Pakistan, Taiwan strait, Falklands)? For the Russian
    provider, whose responses are withheld, the question becomes
    what its refusal behaviour reveals.
  \item[\textbf{RQ3 --- Methodology robustness.}] Is the
    finding an artefact of the OpenAI embedding model used as
    the primary instrument, or does it survive under independent
    embedding spaces (Alibaba, Google Gemini, Yandex)? And is
    the cross-lingual cluster fusion that the embeddings
    recover attributable to subword-tokenizer overlap rather
    than to narrative content?
  \item[\textbf{RQ4 --- What cosine misses.}] What
    language-conditional framing patterns are visible in an
    alternative metric --- a paired-edit stance-axis projection
    --- that anchor cosine flattens? And how does the bias
    generalise when no anchor exists at all, on an imaginary
    contested event for which the model has no training-data
    recall?
\end{description}

\section{Contributions}\label{sec:contributions}

\begin{enumerate}
  \item A reproducible pipeline for measuring linguistic
    ingroup bias in LLM responses against multilingual
    Wikipedia anchors, including a language-axis projection
    debiasing method that separates lexical from framing
    contributions (Chapter~\ref{ch:method}).
  \item A 15-provider $\times$ 9-question $\times$
    3-language sweep on the 2014 Russo-Ukrainian crisis case
    study, spanning seven regulatory regimes; fourteen providers
    across six regimes return analysable prose
    (\S\ref{sec:exp_014}), followed by a five-conflict
    scale-up to Israel--Palestine, India--Pakistan, the Taiwan
    strait, and the Falklands (\S\ref{sec:exp_006}). The
    cosine ingroup-pull pattern generalises across all five
    conflict families, with raw EN ingroup-pull magnitudes
    ranging from ${+}0.32$ (India--Pakistan) down to ${+}0.05$
    (Falklands) --- a raw ordering that anchor-geometry ceilings
    compress substantially (\S\ref{sec:exp_006}). The Falklands near-null functions as a
    falsification anchor for the pipeline --- the metric does not
    return a strong signal everywhere it is pointed --- and the
    event-ladder study of \S\ref{sec:exp_007} rules out event age
    as its explanation, since far older events return far stronger
    signals.
  \item A four-embedder cross-corpus robustness check showing
    that the \emph{sign} of ingroup pull is method-robust but
    the \emph{magnitude ordering} (EN ${\gg}$ RU ${>}$ UK) is
    OpenAI-specific (\S\ref{sec:exp_015}). The Russian-trained
    Yandex embedder gives Ukrainian responses \emph{more}
    ingroup pull than Russian responses, falsifying the simple
    ``Russian-trained embedder will inflate Russian pull''
    hypothesis.
  \item A subword-tokenizer-overlap control ruling out shared
    Cyrillic vocabulary as the dominant explanation for the
    cross-lingual cluster fusion that unsupervised structure
    on the response embeddings recovers
    (\S\ref{sec:exp_017}). All three inspectable embedder
    tokenizers fall below the BLOOM Russian--Ukrainian overlap
    reference of $0.76$ used by \citet{qi2023crosslingual}.
  \item A paired-edit stance-axis projection metric that
    triangulates with anchor cosine and reveals, across the
    same five conflicts, that the Russian--Ukrainian pair shows
    the \emph{smallest} per-language framing gap on this
    measure --- demoting the flagship pair from mid-ranking on
    cosine to the floor. India--Pakistan
    tops the stance ranking with a mean cross-language gap of
    ${+}0.131$; Russian--Ukrainian sits at ${+}0.033$. The
    gap exists in every conflict but its magnitude is
    conflict-specific in a way cosine flattens
    (\S\ref{sec:exp_021}).
  \item Three categorical findings beyond cosine: state-aligned
    refusal in a 270/270 sweep of YandexGPT
    (\S\ref{sec:exp_011}); language-conditional topic
    avoidance in a US-trained diffusion LLM (Mercury-2's
    Crimea coverage drops from a Ukrainian median of $2$ to a
    Russian median of $0$ on the same prompt under a
    pre-specified four-level rubric, \S\ref{sec:exp_003});
    and stylistic-fingerprint differences between frontier
    and 7B local models that are invisible to anchor cosine
    but recoverable from the unsupervised embedding structure
    (\S\ref{sec:exp_016}).
  \item An event-ladder study testing whether the recency gradient
    documented in the human Wikipedia data transfers to model
    responses (\S\ref{sec:exp_007}). Six contested disputes are
    each given a ladder of focal events spanning as much of their
    documented history as Wikipedia supports, and two settled
    separations provide single-event baselines --- 31 events from
    988 to 2023, in sixteen languages, with 31{,}457 completions
    scored on the refusal arm. Ingroup pull is positive on every
    event tested and both settled baselines sit at the floor of the
    contested range in raw pull (below it under length-matched
    anchors, narrowly at the 2022 cell; the floor is partly an
    anchor-geometry effect, \S\ref{sec:exp_007}), but the recency
    gradient does \emph{not}
    transfer: within-family correlations between event date and
    pull split three against three. Bias is inherited in direction,
    not in temporal structure. Four measurement arms --- anchor
    cosine, stance projection, refusal rate, and an LLM judge ---
    agree on direction and diverge on how cleanly they
    separate contested disputes from settled ones.
  \item A two-shape characterisation of bias on imaginary
    contested events. Free-form prose carries content
    divergence without direction; direct attribution carries
    directional commitment without content divergence. On a
    fictional 2024 Russo-Ukrainian bridge incident with
    invented place names absent from training data, the
    free-form news-report question carries the highest
    cross-lingual divergence in the bank with zero net stance
    direction, while the attribution questions carry the
    lowest divergence but the strongest directional lean
    (Ukrainian-ward under Ukrainian prompts); one 7B provider
    shows a 1.07-point swing in LLM-as-judge stance score
    between its Russian-prompt and Ukrainian-prompt means on
    the same fictional event (\S\ref{sec:exp_004}).
\end{enumerate}

\section{Thesis structure}\label{sec:structure}

Chapter~\ref{ch:background} reviews the Wikipedia ingroup-bias
literature, the embedding-space bias measurement methodology
from \citet{bolukbasi2016man} through
\citet{caliskan2017semantics} and \citet{may2019}, the recent
multilingual LLM bias work, and the four published precedents
that position this thesis as a rigorous extension rather than
a fresh discovery. Chapter~\ref{ch:method} describes the
pipeline in detail: prompt design, anchor selection, embedding,
language-axis debiasing, the unsupervised cluster-structure
analysis, and the paired-edit stance-axis projection metric.
Chapter~\ref{ch:experiments} presents the experiments as an
argument rather than in the order they ran: the core finding
first, then three independent attempts to kill it as a
measurement artefact, then the categorical phenomena and
alternative metrics that recover what cosine averages away, then
its generalisation across conflict families, and finally the
event-ladder study that varies time within a conflict. Chapter~\ref{ch:discussion} interprets the
findings, resolves the four-way disagreement between the
measurement instruments, lists threats to validity, and answers
the research questions. Reproduction
recipes, the full prompt bank, the provider matrix, and the cost
summary are in Appendix~\ref{ch:appendix}; the paired-edit
lexicons are frozen in the project repository
(\S\ref{sec:stance_axis}).
